\chapter{Progettazione e codifica}
\label{cap:progettazione-codifica}

\intro{In questo capitolo viene descritta la progettazione 
del sistema e le tecnologie utilizzate per lo sviluppo del prodotto 
software.}\\

\section{Tecnologie e strumenti}
\label{sec:tecnologie-strumenti}
% Di seguito viene data una panoramica delle tecnologie e degli strumenti 
% presi in considerazione durante lo sviluppo del \textit{Proof of Concept}, 
% e successivamente motivare le scelte effettuate.
% Le scelte architetturali hanno preso in considerazione 
Il prodotto software sviluppato ha l'obiettivo di consentire agli utenti 
di interagire con un \textit{chatbot} capace di selezionare, in base alla 
richiesta formulata, il \gls{toolg} più appropriato tra quelli esposti 
dal server \gls{mcpg}, inclusa la funzionalità di \gls{crawlingg} 
implementata in questo progetto. 

La realizzazione di tale sistema richiede una comprensione approfondita 
delle tecniche di \textit{caching} semantico e delle tecnologie di intelligenza 
artificiale applicate al recupero e all'elaborazione delle informazioni testuali.

Le tecnologie destinate al progetto sono state selezionate con l'obiettivo 
di rispondere, nella maniera più efficiente possibile, alle seguenti 
sfide architetturali e funzionali:
\begin{itemize}
    \item \textbf{Gestione dei dati e autenticazione:} amministrazione sicura delle informazioni e profilazione degli utenti;
    \item \textbf{Ottimizzazione del \textit{crawling}:} ottimizzazione del processo di estrazione dei contenuti dalle pagine web;
    \item \textbf{Caching semantico:} progettazione di un meccanismo di memorizzazione intermedia per ridurre le richieste ridondanti all'\gls{llmg}, ottimizzando tempi e costi;
    \item \textbf{Interfaccia utente reattiva:} sviluppo di una \textit{UI} dinamica, capace di mostrare in tempo reale i singoli passaggi del processo di ricerca;
    \item \textbf{Integrazione dell'ecosistema MCP:} implementazione di un server \gls{mcpg} che esponga il servizio di \gls{crawlingg} come strumento standardizzato, rendendolo fluidamente selezionabile dal \textit{chatbot} tra i vari \gls{toolg} disponibili.
\end{itemize}

\subsection{Strumenti utilizzati}
In questa sezione vengono discusse le principali tecnologie 
individuate durante le prime settimane di tirocinio e
utilizzate nello sviluppo del progetto.

\subsubsection{React}

React è una libreria JavaScript \textit{open source} ampiamente utilizzata per la costruzione di interfacce utente dichiarative e basate su componenti. 
La sua adozione in questo progetto riflette la necessità di costruire un'interfaccia complessa, reattiva e facilmente manutenibile per la gestione delle conversazioni e degli utenti.
React si basa sul concetto di componente: ogni elemento dell'interfaccia viene scomposto in unità indipendenti, riutilizzabili e componibili tra loro. 
Questo approccio facilita la suddivisione del lavoro di sviluppo e la manutenzione del codice nel tempo, poiché ogni componente può essere modificato o testato in isolamento 
senza impattare il resto dell'applicazione.
Un altro aspetto rilevante di React è il suo \textit{Virtual DOM}, un meccanismo che consente di aggiornare dinamicamente gli oggetti nelle pagine, 
evitando operazioni di \textit{refresh} dell'intera pagina. Questa caratteristica risulta particolarmente utile in questo progetto, 
dove l'interfaccia deve aggiornarsi di continuo per mostrare in tempo reale i singoli passaggi durante l'elaborazione della risposta 
da parte del \gls{llmg}.
React vanta inoltre un insieme vasto di librerie di supporto, tra cui \textit{React Router} per la gestione della navigazione e \textit{i18next} per l'internazionalizzazione 
dei contenuti, entrambe impiegate in questo progetto. 

Per questi motivi, React rappresenta una scelta solida ed efficace per lo sviluppo del \textit{frontend} di questa applicazione.

\subsubsection{Vite}
Vite è uno strumento di compilazione per applicazioni web moderne, progettato per offrire un'esperienza di sviluppo estremamente rapida. 
Infatti, Vite sfrutta i moduli \textit{ES} (\textit{ECMAScript Modules}) nativi del \textit{browser} durante la fase di sviluppo, 
evitando di dover ricompilare l'intera applicazione ad ogni modifica del codice.
Questo strumento si distingue per i tempi di avvio quasi istantanei 
e permette di visualizzare le modifiche al codice 
nel \textit{browser} in tempi minimi, senza dover ricaricare l'intera pagina.

La scelta di Vite integrata con React è motivata dalla necessità di mantenere un flusso di sviluppo rapido 
ed efficiente durante l'iterazione continua sui componenti dell'interfaccia.

\subsubsection{Python}
Python è un linguaggio di programmazione ad alto livello, interpretato e di uso generale, noto per la sua sintassi chiara e leggibile. 
Tale caratteristica favisce la leggibilità del codice e riduce i costi di manutenzione del progetto, 
aspetto particolarmente rilevante in un'architettura suddivisa in molti microservizi indipendenti.
Python si distingue per la sua versatilità: può essere impiegato per lo sviluppo \textit{backend}, 
l'elaborazione di dati, l'integrazione con modelli di intelligenza artificiale e l'automazione di processi. 
Questa flessibilità lo rende particolarmente adatto a un progetto come questo, in cui il linguaggio deve gestire 
sia la logica di business dei vari servizi sia l'orchestrazione delle chiamate al \gls{llmg}. 

Per questi motivi, Python rappresenta la scelta naturale per l'implementazione di tutti i servizi \textit{backend} del sistema.

\subsubsection{FastAPI}

FastAPI è un \gls{frameworkg} web moderno per Python, progettato per la costruzione di \gls{apig} ad alte prestazioni. 
Si basa sugli standard \textit{OpenAI} 
e \textit{JSON Schema}, e sfrutta i \textit{type hints} di Python per la validazione automatica dei dati in ingresso e in uscita.
Una delle caratteristiche principali di FastAPI è la generazione automatica della documentazione interattiva delle API, 
basata sugli schemi \textit{Pydantic} definiti nel codice. 
Questo riduce sensibilmente il tempo necessario per la redazione della specifica degli \textit{endpoint} e garantisce che 
quest'ultima sia sempre sincronizzata con l'implementazione effettiva.
FastAPI è inoltre costruito su \textit{Starlette} e supporta nativamente la programmazione asincrona, 
aspetto fondamentale per questo progetto, dove i singoli microservizi devono gestire un elevato numero di richieste e comunicazioni in tempo reale 
tramite \gls{sseg}. 

Per queste ragioni, FastAPI è stato scelto come \gls{frameworkg} principale per l'implementazione dei servizi 
\textit{backend} del sistema, secondo un'architettura a porte e adattatori (\textit{hexagonal architecture}).

\subsubsection{FastMCP e il Model Context Protocol}

Il \gls{mcpg} è un protocollo standardizzato che consente ai modelli di intelligenza artificiale di interagire 
con strumenti e servizi esterni tramite un'interfaccia uniforme, eliminando la necessità di implementare 
integrazioni specifiche per ciascun provider o strumento. 
FastMCP è un \gls{frameworkg} Python che semplifica notevolmente l'implementazione di un server 
\gls{mcpg}, occupandosi della gestione del protocollo, della definizione degli strumenti esposti e della comunicazione con i client.
FastMCP si distingue per la sua semplicità d'uso: permette di esporre funzioni Python come strumenti 
utilizzabili da un modello di linguaggio tramite semplici decoratori (come \texttt{@mcp.tool}), senza richiedere la gestione 
manuale dei dettagli di basso livello del protocollo.

Nel contesto di questo progetto, FastMCP rappresenta il cuore del sistema: il server espone, 
tra i propri strumenti, il \gls{toolg} di \gls{crawlingg}, che il \textit{chatbot} è in grado di selezionare 
automaticamente quando il messaggio dell'utente lo richiede. 
La scelta di FastMCP consente quindi di concentrarsi sulla logica di estrazione e ricerca dei contenuti, 
delegando a tale strumento la gestione del protocollo di comunicazione e di esposizione degli strumenti al modello.

\subsubsection{BeautifulSoup}
BeautifulSoup è una libreria Python per l'estrazione del contenuto da documenti \gls{htmlg} e \gls{xmlg}\glsfirstoccur .
Questa libreria si distingue per la sua tolleranza nei confronti di \gls{htmlg} malformato, una caratteristica particolarmente 
utile quando si analizzano pagine web reali, che spesso non rispettano rigorosamente gli standard. 
BeautifulSoup fornisce inoltre un insieme di metodi per la ricerca e il filtraggio degli elementi della pagina in base 
a \textit{tag}, classi, identificatori o altri attributi, semplificando notevolmente l'estrazione delle informazioni rilevanti.

Nel contesto di questo progetto, BeautifulSoup viene impiegata nella fase di recupero delle pagine web permettendo di
prelevare il loro contenuto e i collegamenti ipertestuali, 
necessari per la generazione dei riassunti e la scelta 
dei \textit{link} per il processo di \gls{crawlingg}.

\subsubsection{OpenAI}
OpenAI è un'organizzazione di ricerca e sviluppo nel campo 
dell'intelligenza artificiale, celebre i modelli 
GPT (\textit{Generative Pre-trained Transformer}).
I modelli offerti da OpenAI sono accessibili tramite servizi \textit{cloud} e 
possono essere utilizzati attraverso \gls{apig}, permettendo di 
realizzare sistemi basati sulla comprensione del linguaggio naturale.

Nel contesto di questo progetto, OpenAI rappresenta il componente 
centrale per l'elaborazione delle richieste 
degli utenti e la generazione delle risposte.

\subsubsection{LangSmith}

LangSmith è una piattaforma di tracciabilità e \textit{debugging} pensata specificamente per applicazioni basate 
su modelli di intelligenza artificiale. 
Permette di tracciare in modo dettagliato l'esecuzione delle catene di chiamate che coinvolgono \gls{llmg}, 
fornendo una visione completa del flusso di esecuzione.
Questa piattaforma si distingue per la sua capacità di raccogliere \textit{trace} distribuite anche in architetture 
a microservizi, dove una singola richiesta dell'utente può attraversare più servizi prima di giungere a 
una risposta finale. 
Ogni passaggio viene registrato con i relativi \textit{input}, \textit{output}, tempi di esecuzione e metadati, 
permettendo di individuare rapidamente latenze, errori o comportamenti inattesi del modello.

Nel contesto di questo progetto, LangSmith è stato adottato per il tracciamento distribuito delle operazioni 
svolte durante il processo di \gls{crawlingg}, dal momento della selezione del \gls{toolg} da parte 
del \textit{chatbot} fino alla generazione della risposta finale. 
Questo strumento si è rivelato fondamentale durante le fasi di sviluppo e \textit{debug} del sistema, 
permettendo di analizzare nel dettaglio il comportamento del modello.

\subsubsection{PostgreSQL, TimescaleDB e pgvector}
PostgreSQL è un sistema di gestione di basi di dati relazionali \textit{open source}, noto per la sua robustezza, 
l'affidabilità e la conformità agli standard \gls{sqlg}\glsfirstoccur . 
La sua espandibilità tramite estensioni native lo rende particolarmente adatto a progetti che richiedono funzionalità 
avanzate oltre alla semplice gestione di dati relazionali, come nel caso di questo progetto.
TimescaleDB è un'estensione di PostgreSQL pensata per la gestione efficiente di dati temporali (\textit{time-series}), 
ottimizzata per l'inserimento e l'interrogazione di grandi volumi di dati ordinati nel tempo. 
Tale estensione è stata adottata per la gestione dello storico delle ricerche effettuate dagli utenti, 
dove ogni elemento è naturalmente associato a una dimensione temporale.
L'estensione pgvector introduce invece un tipo di dato vettoriale e operatori di similarità, 
permettendo di memorizzare e interrogare \textit{embedding} direttamente all'interno del database relazionale. 
Questa estensione è fondamentale per l'implementazione del meccanismo di \textit{caching} semantico del sistema: 
le domande poste dagli utenti vengono trasformate in vettori numerici e confrontate con quelle già elaborate in precedenza, 
individuando richieste semanticamente simili ed evitando così rielaborazioni ridondanti da parte dell'\gls{llmg}.

La combinazione di queste tre tecnologie permette di utilizzare un unico database per gestire sia 
i dati relazionali tradizionali (utenti, storico delle conversazioni, profili) sia i dati vettoriali necessari al 
\textit{caching} semantico, semplificando l'infrastruttura complessiva del sistema.

\subsubsection{Docker e Docker Compose}
Docker è una piattaforma di containerizzazione che consente di raggruppare un'applicazione e tutte le sue dipendenze 
in un'unità isolata e portabile, denominata ''container''. 
Questo approccio garantisce che l'applicazione si comporti nello stesso modo indipendentemente dall'ambiente 
in cui viene eseguita, eliminando i classici problemi di incompatibilità tra ambienti di sviluppo, test e produzione.
Docker Compose è uno strumento che permette di definire e orchestrare applicazioni multi-container 
tramite un singolo file di configurazione dichiarativo. 

Questo strumento risulta particolarmente utile in questo progetto dove l'architettura del sistema è 
suddivisa in molteplici microservizi indipendenti e si richiede scalabilità e portabilità.

\subsubsection{Nginx}
Nginx è un \textit{reverse proxy} \textit{open source}, noto per le sue elevate prestazioni e per la capacità 
di gestire un numero elevato di connessioni simultanee con un basso utilizzo di risorse. 
In questo progetto, Nginx svolge il ruolo di \textit{gateway} per l'intera applicazione, instradando le richieste 
provenienti dal \textit{frontend} verso i corretti microservizi \textit{backend}.
Questo approccio centralizza la gestione del traffico in ingresso, semplificando la configurazione dell'instradamento 
delle richieste e permettendo di disaccoppiare il \textit{frontend} dalla conoscenza della struttura interna 
del sistema a microservizi.

La scelta di Nginx come \textit{gateway} riflette quindi l'esigenza di un punto di ingresso unico, 
performante e configurabile per l'intera architettura distribuita del sistema.

\subsubsection{JWT e JWKS}
\gls{jwtg} è uno standard aperto per la creazione di \textit{token} di accesso che permettono di trasmettere informazioni 
in modo sicuro tra due parti, sotto forma di oggetto \gls{jsong} firmato digitalmente. 
Nel contesto di questo progetto, i \gls{jwtg} vengono utilizzati per gestire l'autenticazione degli utenti tra 
i diversi microservizi del sistema.

In particolare, è stato adottato l'algoritmo di firma \textit{RS256}, basato su una coppia di chiavi asimmetriche: 
una chiave privata, utilizzata dal servizio responsabile dell'autenticazione, per firmare i \textit{token} al momento del \textit{login}, 
e una chiave pubblica, distribuita agli altri servizi tramite \gls{jwksg}\glsfirstoccur , per la verifica della validità 
dei \textit{token}. 

L'adozione di \textit{RS256} e \gls{jwksg} risponde quindi a un'esigenza di sicurezza e di efficienza architetturale, 
garantendo che ogni servizio possa operare in autonomia mantenendo al contempo un meccanismo di autenticazione robusto 
e centralizzato nella sua gestione delle chiavi.

% \section{Ciclo di vita del software}
% \label{sec:ciclo-vita-software}

\section{Progettazione}
\label{sec:progettazione-architetturale}
L'infrastruttura è stata definita secondo un modello stratificato (\textit{multitier}), 
strutturato su due livelli macroscopici: il livello di presentazione 
(\textit{frontend}) e lo strato di logica applicativa (\textit{backend}).

\subsection{Progettazione backend} 
Il \textit{backend} rappresenta il nucleo computazionale dell'applicazione. 
Viene suddiviso secondo l'architettura a microservizi, 
una scelta mirata a garantire un elevato livello di modularità e di 
separazione delle responsabilità.

L'intero livello è sviluppato in linguaggio Python 
sfruttando i \gls{frameworkg} FastAPI e FastMCP. 
A seguire vengono descritti nel dettaglio i singoli microservizi 
e le principali sfide affrontate nel loro sviluppo.

\subsubsection{MCP Server Service}
Il MCP Server Service si basa sul protocollo \gls{mcpg} ed è 
responsabile dell'esposizione del \gls{toolg} di \textit{crawling}. 
Esso si occupa, inoltre, di generare gli eventi \gls{sseg} e inviarli 
verso il \textit{frontend} per notificare all'utente, in tempo reale, 
i singoli passaggi logici del processo.

La difficoltà maggiore durante la codifica di tale servizio ha riguardato 
l'ingegnerizzazione della logica di ricorsione del processo di \textit{crawling}. 
Ad ogni passo ricorsivo, l'algoritmo esegue in modo sequenziale i seguenti passaggi:
\begin{enumerate}
    \item \textbf{Estrazione:} recupero del contenuto \gls{htmlg} della pagina corrente $X$;
    \item \textbf{Sintesi:} generazione di un riassunto mirato in base alla richiesta originaria dell'utente, il quale viene accodato alla lista dei riassunti raccolti nei passi precedenti;
    \item \textbf{Valutazione del contesto:} verifica della sufficienza informativa dei dati raccolti per la generazione della risposta, effettuata analizzando l'intera lista dei riassunti accumulati. Tale controllo determina due possibili scenari:
    \begin{itemize}
        \item \textbf{Successo}: il contesto è valutato come esaustivo per rispondere al quesito; la ricorsione viene interrotta e si procede alla generazione della risposta finale;
        \item \textbf{Prosecuzione}: i riassunti correnti sono considerati insufficienti per formulare una risposta completa, rendendo necessaria la continuazione del processo di \textit{crawl};
    \end{itemize}
    \item \textbf{Parsing dei collegamenti:} in caso di mancata risoluzione al punto precedente, vengono estratti dalla pagina $X$ tutti i collegamenti ipertestuali disponibili;
    \item \textbf{Filtraggio semantico:} selezione dei \textit{link} ritenuti maggiormente pertinenti rispetto alla domanda dell'utente;
    \item \textbf{Passo ricorsivo:} per ciascun collegamento selezionato, viene avviata una nuova istanza della medesima funzione ricorsiva.
\end{enumerate}
La logica sopra descritta modella l'esplorazione secondo una struttura ad albero. 
In questa analogia algoritmica, l'\textbf{ampiezza} della struttura 
è determinata dal fattore di \textit{branching}, ossia dal numero massimo di collegamenti 
selezionati durante la fase di filtraggio semantico (passo 5); l'\textbf{altezza} 
dell'albero (ovvero la sua profondità massima) corrisponde invece al limite massimo 
di chiamate ricorsive nidificate lungo il cammino più lungo che intercorre tra il 
nodo radice e un nodo foglia.


\paragraph{Complessità ricorsiva}
La principale criticità dell'algoritmo descritto risiede nella 
definizione dei parametri di ampiezza e profondità. La configurazione di tali valori
richiede una calibrazione 
estremamente rigorosa: una scelta sovradimensionata esporrebbe il sistema al 
fenomeno dell'\textit{esplosione combinatoria}. 
Ciò comporterebbe un consumo eccessivo di risorse computazionali e un incremento 
incontrollato delle chiamate verso l'\gls{llmg}, con conseguenti ripercussioni in 
termini di latenza e costi finanziari. 
Di contro, l'adozione di vincoli stringenti 
comprometterebbe l'efficacia stessa del processo 
di ricerca: in caso di domande più complesse, dei parametri troppo bassi condurrebbero 
inevitabilmente a un'analisi superficiale e, di conseguenza, a un 
fallimento nel soddisfare la richiesta dell'utente.
\\\\La Tabella~\ref{tab:analisi-casi-ricorsione} mostra
come i parametri di ampiezza e profondità 
influenzano la complessità ricorsiva
nel caso peggiore (ovvero qualora l'algoritmo esaurisse 
l'intero albero senza raggiungere prima lo stato di successo). 
Nel processo di \textit{crawling}, ogni singolo passo (nodo dell'albero) 
richiede esattamente tre chiamate distinte all'\gls{llmg} (fasi di sintesi, 
valutazione del contesto e selezione dei \textit{link}).
\begingroup
\renewcommand{\arraystretch}{1.5}
\begin{xltabular}{\textwidth}{|c|c|c|c|c|}
\hline
\textbf{Caso} & \textbf{Ampiezza} & \textbf{Profondità} & \textbf{\makecell{Chiamate di ricorsione}} & \textbf{\makecell{Chiamate al LLM}} \\
\hline
1 & 3 & 6 & 1092 & 3276 \\
\hline
2 & 1 & 4 & 4 & 12 \\
\hline
3 & 2 & 6 & 126 & 378 \\
\hline
4 & 2 & 5 & 62 & 186 \\
\hline
\caption{Analisi dell'impatto dei parametri di profondità e ampiezza sul numero di chiamate.}
\label{tab:analisi-casi-ricorsione}
\end{xltabular}
\endgroup

Sulla base delle analisi effettuate durante la fase di \textit{testing} e bilanciando il \textit{trade-off} tra l'accuratezza del reperimento informativo (\textit{recall}) e il contenimento di costi e latenze, la configurazione ottimale è stata individuata nei seguenti valori:
\begin{itemize}
    \item \textbf{Ampiezza = 2}: Limitare il fattore di \textit{branching} a un massimo di due collegamenti ipertestuali per pagina previene l'espansione iperbolica dell'albero;
    \item \textbf{Profondità = 5}: Configurare l'altezza massima a cinque livelli di nidificazione consente all'algoritmo di penetrare sufficientemente a fondo nella struttura del sito web.
\end{itemize}

Con questa specifica impostazione ($b=2, d=5$), il numero massimo teorico di nodi 
esplorabili è pari a $N_{\max} = 2^1 + 2^2 + 2^3 + 2^4 + 2^5 = 62$. 
Di conseguenza, sempre nel caso peggiore, il numero massimo di chiamate all'\gls{llmg} si attesta a:
\begin{equation}
    C_{tot} = 62 \times 3 = 186 \text{ chiamate}
\end{equation}
Tale valore rappresenta un limite sostenibile per l'infrastruttura 
e perfettamente in linea con le capacità operative dei modelli di intelligenza 
artificiale aziendali.

\paragraph{Collegamenti ipertestuali ciclici}
Un'ulteriore criticità strutturale è rappresentata dalla presenza di 
collegamenti ipertestuali ciclici (o \textit{link} circolari) all'interno 
delle pagine web esplorate. 
In assenza di un opportuno meccanismo di tracciamento, l'algoritmo 
rischierebbe di rimanere intrappolato in cicli di navigazione infiniti tra 
pagine che si referenziano reciprocamente.

Per ovviare a questo problema, la soluzione più efficiente e lineare prevede 
l'introduzione di una struttura dati per la memorizzazione degli stati visitati, 
definita a livello di codice come una lista denominata \texttt{visited}. 
All'avvio di ogni passo ricorsivo, la funzione registra l'\gls{urlg} della pagina 
corrente all'interno di tale struttura, trasmettendola poi come parametro alle 
successive chiamate. 
Prima di procedere alla fase di scelte dei collegamenti (passo 5), 
l'algoritmo esegue un'operazione di sottrazione insiemistica, 
rimuovendo dall'elenco dei \textit{link} candidati tutti gli \gls{urlg} già 
presenti in \texttt{visited}.

\paragraph{Caching semantico}
A differenza delle tecniche di memoizzazione tradizionali 
(come il \textit{caching} basato su chiavi e valori esatti o \textit{hash} 
di stringhe), il \textit{caching} semantico valuta il significato logico 
e contestuale delle richieste inviate dagli utenti. 
Questa progettazione risponde all'esigenza di ottimizzare 
il sistema evitando di rielaborare richieste rindondanti e, 
di conseguenza, consumare inutilmente \textit{token} del \gls{llmg}.

Il sottosistema si appoggia perciò sull'estensione \textit{pgvector} di PostgreSQL e opera 
secondo una logica sequenziale:
\begin{enumerate}
    \item \textbf{Generazione dell'embedding}: il messaggio dell'utente viene convertito 
    in un vettore chiamato \textit{embedding} ad alte dimensioni, 
    capace di catturarne i tratti semantici;
    \item \textbf{Ricerca di similarità}: il vettore generato viene utilizzato per interrogare 
    la tabella di \textit{cache} sul database relazionale.
    Attraverso la funzione di calcolo di similarità del coseno, il sistema 
    confronta l'\textit{input} corrente con le interrogazioni precedentemente storicizzate 
    nelle ultime ore;
    \item \textbf{Valutazione della soglia e raffinamento:} vengono estratti dal database i 
    record che presentano la massima similarità riscontrata. 
    Successivamente, per garantire la massima accuratezza ed evitare falsi positivi, 
    viene selezionato il record semanticamente più pertinente tramite una chiamata di 
    validazione al \gls{llmg}. In base all'esito di questo controllo si determinano due scenari:
    \begin{itemize}
        \item \textbf{Cache Hit}: il modello conferma l'equivalenza semantica; 
        il sistema estrae la risposta pre-calcolata e la restituisce, 
        ottimizzando tempi e costi;
        \item \textbf{Cache Miss}: non viene riscontrata alcuna corrispondenza; 
        viene avviata la normale elaborazione della richiesta.
    \end{itemize}
\end{enumerate}
Tale strategia di memoizzazione trasforma l'infrastruttura da un sistema puramente reattivo 
e computazionalmente oneroso ad uno in grado di incrementare l'efficienza 
e la qualità del prodotto software.

\paragraph{Invio delle notifiche sui passaggi di crawling}
Un'altra difficoltà affrontata è legata alla gestione degli eventi 
da inviare al \textit{frontend} riguardanti i passaggi logici del \textit{crawler}. 
Poiché l'esplorazione ricorsiva e l'elaborazione dei contenuti richiedono tempistiche 
computazionali non immediate, l'utilizzo di una classica architettura \gls{httpg} 
di tipo sincrono (\textit{request-response}) avrebbe esposto il sistema a frequenti 
fenomeni di \textit{timeout}, oltre a lasciare l'utente in uno stato di attesa privo 
di \textit{feedback} visivo.
La soluzione ha dunque previsto l'implementazione di un modello di 
comunicazione asincrono basato su \gls{sseg}. 
Il MCP Server Service espone un \textit{endpoint} per la definizione di un 
canale di streaming con il \textit{frontend}, consentendo di notificare 
in tempo reale ogni singolo stato di avanzamento della ricerca 
("\textit{Prelevata pagina $X$}", "\textit{Generazione sintesi}", 
"\textit{Scelta dei link}").


\subsubsection{Search Service}
Il servizio Search Service è dedicato alla ricezione e 
all'elaborazione dei messaggi dell'utente espressi in linguaggio naturale. 
Analizzando il contesto della richiesta, il servizio determina quale 
\gls{toolg} sia necessario attivare, ne estrae i parametri 
e invoca l'\textit{endpoint} corrispondente sul microservizio 
MCP Server Service. 
Nonostante i requisiti del progetto richiedessero esclusivamente lo 
sviluppo del \gls{toolg} di \textit{crawling}, il componente 
è stato nativamente concepito per poter gestire ulteriori 
\gls{toolg} in futuro; l'esistenza del \textit{search service} 
come orchestratore e selettore dinamico delle azioni risponde proprio 
a questo fondamentale vincolo di estensibilità.

\paragraph{Iniezione della cronologia e riordino dei messaggi}
Una delle principali sfide riguarda la corretta ricostruzione della conversazione: 
per consentire all'\gls{llmg} di comprendere appieno le richieste dell'utente, 
il Search Service non elabora il singolo messaggio isolato, bensì recupera 
l'intero set di messaggi presenti nella \textit{chat}. 
Questa sequenza di messaggi viene strutturata e ordinata in senso cronologico 
decrescente (dal più recente al più vecchio): tale gerarchia informativa è 
fondamentale per specificare al modello quale sia il messaggio principale
e l'obiettivo dell'utente, evitando che il contesto pregresso comprometta 
la coerenza della risposta.

\paragraph{Estrazione dei parametri e non-determinismo}
Un'altra difficoltà riguarda il non determinismo dei modelli \gls{llmg}. 
Quando l'utente formula le richieste in linguaggio naturale, 
il \textit{search service} deve garantire che l'\gls{llmg} non solo selezioni 
lo strumento corretto, ma estragga i parametri di input (come \gls{urlg}, parole 
chiave o vincoli di ricerca) in un formato rigidamente tipizzato e privo di ambiguità.

Per risolvere tale criticità, la progettazione ha previsto l'adozione di schemi di 
validazione dei dati basati sulla libreria \textit{Pydantic}. 
Forzando il modello a rispondere secondo uno schema \gls{jsong} rigido, 
si è ridotto quasi allo zero il rischio di ricevere risposte malformate in 
caso di "allucinazioni".

\subsubsection{Auth Service}
L'Auth Service gestisce l'intero ciclo di vita dell'autenticazione 
e della sicurezza degli accessi. 
Ha la responsabilità di creare, validare e revocare i \gls{jwtg} 
(gestendo sia gli \textit{access token} che i \textit{refresh token}). 
Per la firma dei \textit{token} si è scelto di adottare l'algoritmo asimmetrico 
\textit{RS256} (\textit{RSA} utilizzando \textit{SHA-256}): 
questo approccio garantisce che solo l'\textit{auth-service}, custode 
della chiave privata, possa generare e firmare \textit{token} validi, 
consentendo al contempo agli altri microservizi di verificarne l'autenticità 
in totale autonomia tramite la rispettiva chiave pubblica, senza dover 
effettuare chiamate inter-servizio sincrone per ogni richiesta. 
A tale scopo, il servizio espone un \textit{endpoint} chiamato 
''\texttt{.well-known/jwks.json}'' che permette di distribuire la chiave pubblica 
tra i vari microservizi del\textit{backend}.

\subsubsection{User Service}
User Service è il servizio dedicato alla creazione e alla gestione degli 
utenti della piattaforma. Offre le seguenti funzionalità principali:
\begin{itemize}
    \item \textbf{Creazione utente}: crea un nuovo utente dato username, email e 
    password. Solo gli amministratori possono creare un utente;
    \item \textbf{Modifica dati utente}: un utente può modificare 
    username o password accedendo al proprio profilo;
    \item \textbf{Modifica permessi ai tool di un utente}: un admin può 
    modificare i vari permessi di accesso, utilizzo della \textit{cache} e 
    visualizzazione dei passaggi di ricerca di un utente base;
    \item \textbf{Eliminazione utente}: un admin può eliminare un utente 
    con qualsiasi ruolo;
\end{itemize}


\subsubsection{History Service}
History Service si occupa della persistenza e della sincronizzazione dello storico 
delle conversazioni, memorizzando le conversazioni tra l'utente e l'\gls{llmg} 
per consentire il recupero e la consultazione delle \textit{chat} passate. 
Oltre a garantire la continuità dell'esperienza utente tra i diversi accessi, 
questo componente svolge un ruolo chiave nel fornire al modello il contesto 
storico necessario per comprendere la richiesta dell'utente all'interno della 
stessa \textit{chat}.

\paragraph{Efficienza della persistenza e disaccoppiamento guidato dal client}
Un problema riguarda l'impatto prestazionale derivante dalle frequenti 
operazioni di scrittura sul database relazionale.
Per evitare che il salvataggio della 
cronologia generi latenza o blocchi il flusso principale 
di generazione della risposta, l'architettura sfrutta un disaccoppiamento 
basato sull'orchestrazione da parte del \textit{frontend}. 
Il client, infatti, invoca inizialmente il Search Service per avviare il 
processo di ricerca e, solo al suo termine, effettua 
una chiamata verso l'\textit{endpoint} dell'\textit{History Service} per 
salvare i nuovi messaggi. 
Questo approccio garantisce che i tempi di scrittura 
sul database non interferiscano mai con la generazione della risposta.


% \begin{namespacedesc}
%     \classdesc{Classe 1}{Descrizione classe 1}
%     \classdesc{Classe 2}{Descrizione classe 2}
% \end{namespacedesc}


\subsection{Progettazione frontend}
Il \textit{frontend} rappresenta il livello di presentazione dell'applicazione, 
ovvero l'interfaccia diretta tra l'utente finale e i microservizi del \textit{backend}. 
La reattività e la fluidità dell'interfaccia grafica costituiscono requisiti 
non funzionali fondamentali per garantire un'esperienza d'uso equiparabile a 
quella dei moderni sistemi di \textit{chatbot}.

Poiché la pagina di \textit{chat} deve mostrare i passaggi 
logici di ricerca in tempo reale (come il recupero delle pagine, 
la sintesi del contenuto della pagina e il filtraggio dei collegamenti) 
senza comportare il ricaricamento di tutta la pagina, si è scelto di adottare 
il pattern architetturale \gls{mvvmg} (Model-View-ViewModel). 
Questa separazione delle responsabilità permette di disaccoppiare 
rigidamente l'interfaccia utente (\textit{View}) dalla logica di business 
e dallo stato dei dati (\textit{Model}). 
Il \textit{ViewModel} agisce come intermediario: quando un determinato oggetto 
cambia valore, aggiorna automaticamente lo stato della vista, garantendo una resa 
grafica fluida.

\paragraph{Controllo degli accessi e gestione delle rotte}
Al fine di garantire la sicurezza dell'applicazione e una corretta separazione 
dei permessi a livello di interfaccia utente, la logica di navigazione del 
\textit{frontend} implementa un sistema di protezione delle rotte basato su 
componenti \textit{wrapper} (componenti di ordine superiore) sfruttando la 
libreria \textit{React Router}. 
L'accesso alle diverse macro-aree dell'applicazione (autenticazione, 
chat e pannello di amministrazione) viene regolato controllando 
lo stato del profilo utente mediante chiamate asincrone all'\textit{endpoint} di 
sessione (\texttt{getCurrentUser}). 
Tale strategia garantisce che un utente del sistema possa visualizzare e 
interagire solamente con le viste a cui ha accesso.
I moduli di controllo sono \texttt{Public Route}, \texttt{ProtectedRoute} e 
\texttt{AdminRoute}.

Il componente \texttt{Public Route} gestisce l'unica pagina accessibile da utenti non 
autenticati, ovvero la pagina di \textit{login}. 
Il problema risolto da questo modulo risiede nell'impedire a un utente già 
autenticato di accedere nuovamente a moduli di login superflui. 
Qualora la verifica restituisca uno stato \texttt{'authenticated'}, il 
componente esegue un reindirizzamento forzato verso pagina principale di \textit{chat}.

Al contrario, se un utente non autorizzato tenta di accedere ad una pagina diversa 
da quella di \textit{login}, viene automaticamente reindirizzato a quest'ultima.
L'utilizzo della proprietà \texttt{replace} è fondamentale in termini di gestione 
della cronologia del browser (\textit{history stack}), poiché sostituisce 
l'indirizzo corrente impedendo all'utente di tornare indietro alla pagina 
non consentita tramite i comandi nativi del browser.

Per l'accesso alle principali funzionalità come \textit{chat}, gestione utenti 
e profilo utente, sono stati sviluppati i moduli 
\texttt{ProtectedRoute} e \texttt{AdminRoute}. 
Il primo si assicura che il \textit{token} di sessione sia valido: in caso di errore o di 
sessione scaduta (\texttt{'expired'}), l'utente viene reindirizzato alla 
pagina di \textit{login}. Il componente \texttt{AdminRoute} estende questa 
logica introducendo un meccanismo di segregazione dei ruoli: dopo aver 
accertato l'autenticazione, il componente verifica che il ruolo dell'utente 
sia ''admin''. 
Se l'utente non possiede i privilegi di amministratore, avviene un reindirizzamento 
verso la pagina della \textit{chat}. 
\\\\Questa gerarchia strutturata garantisce che ogni 
tipologia di utente si trovi all'interno del proprio perimetro 
operativo.


\section{Design Pattern utilizzati}
L'architettura del sistema è stata strutturata adottando molteplici 
\textit{design pattern} consolidati dall'ingegneria del software. 
L'impiego di questi modelli architetturali risponde alla necessità di 
garantire il disaccoppiamento tra i componenti e l'estensibilità del 
software.

\subsection{Pattern Singleton}
Il pattern \textit{Singleton} è stato adottato per la gestione della classe 
\texttt{JwtValidation}, responsabile della validazione dei token \gls{jwtg} all'interno 
del \textit{backend}.
L'obiettivo era garantire che esistesse un'unica istanza di tale classe per tutto il
ciclo di vita dell'applicazione, in modo da condividere la \textit{cache} della chiave
pubblica (\texttt{\_jwks\_cache}) tra tutte le richieste ed evitare chiamate
ridondanti all'\textit{endpoint} \texttt{/.well-known/jwks.json} esposto.

Il modulo \texttt{JwtValidation} definisce la propria logica di validazione; la sua istanza, 
invece, viene gestita tramite una metaclasse dedicata, \texttt{SingletonMeta}, in 
maniera tale da rendere il comportamento di \textit{singleton} riutilizzabile e senza duplicare 
la logica di sincronizzazione in ciascuna di esse.

La metaclasse \texttt{SingletonMeta} mantiene un dizionario \texttt{\_instances},
che associa ciascuna classe che la utilizza alla propria istanza unica, e un
oggetto \texttt{\_lock} per garantire la sicurezza in ambiente \textit{multi-thread} durante
la creazione. 
Il metodo \texttt{\_\_call\_\_} viene invocato automaticamente da
Python ogni volta che la classe \texttt{JwtValidation} viene istanziata: se per
quella classe non esiste ancora un'istanza in \texttt{\_instances}, ne viene creata
una sotto lock, altrimenti viene restituita quella già presente.

\begin{figure}[H]
    \centering
    \includegraphics[width=0.8\textwidth]{uml/singleton.png}
    \caption{Diagramma UML del pattern Singleton implementato per \texttt{JwtValidation}}
\end{figure}

Grazie a questa architettura strutturata, il componente \texttt{JwtValidation} garantisce 
un unico punto di accesso globale, memorizza rigidamente in \textit{cache} la chiave 
pubblica dopo il primo recupero e restituisce direttamente l'oggetto validato 
\texttt{JwtClaims} (correlato con il rispettivo set di \texttt{ToolPermission}), 
minimizzando l'\textit{overhead} computazionale e i tempi di risposta 
complessivi dell'infrastruttura.


\subsection{Pattern Adapter}
Il pattern \textit{Adapter} costituisce il meccanismo di implementazione
dell'architettura esagonale (\textit{ports and adapters}) adottata in modo
sistematico in tutti i microservizi del \textit{backend}. Ogni \textit{port}
definito nello strato applicativo (ad esempio \texttt{CachePort},
\texttt{SummarizePort}) descrive un'interfaccia astratta orientata al dominio,
priva di qualunque riferimento a tecnologie o librerie specifiche. Il relativo
\textit{adapter} fornisce invece l'implementazione concreta di tale interfaccia,
traducendo le chiamate provenienti dai casi d'uso in operazioni effettive verso
sistemi esterni (database, chiamate al \gls{llmg}, librerie). Questa separazione
permette di sostituire o modificare l'infrastruttura sottostante senza impattare
la logica applicativa.

A titolo esemplificativo si riportano due \textit{adapter} implementati nel
microservizio MCP Server Service: \texttt{CacheAdapter} e
\texttt{SummarizeAdapter}.

La classe \texttt{CacheAdapter} implementa \texttt{CachePort} e centralizza
tutte le operazioni di \textit{caching} del microservizio. 
L'\textit{adapter} in questione 
realizza il processo di \textit{semantic caching} descritta nel capitolo precedente: 
il metodo \texttt{findCachedRequest}
genera l'\textit{embedding} della domanda tramite \texttt{OpenAIRepository}, recupera i
candidati semanticamente più simili dal repository tramite ricerca per
similarità vettoriale e infine invoca il metodo privato
\texttt{\_getSimilarQuestionIndex}, che sottopone i candidati a un modello
linguistico con un \textit{prompt} di sistema volutamente restrittivo: la
corrispondenza viene accettata solo se il contenuto candidato copre
\emph{esattamente} tutte le informazioni richieste, senza inferenze o
approssimazioni. 
Solo in caso di esito positivo viene restituito un
\texttt{ToolResult} con la risposta in cache; in caso contrario il metodo
restituisce \texttt{None}, segnalando un \textit{cache miss} e lasciando che il
caso d'uso proceda con l'elaborazione della risposta. La stessa logica di
ricerca semantica viene utilizzata da \texttt{findCachedSummary} per la cache
dei riassunti delle pagine.

La classe \texttt{SummarizeAdapter} implementa \texttt{SummarizePort} e adatta
il bisogno applicativo di sintesi del contenuto verso il modello linguistico
esposto da \texttt{OpenAIRepository}. Il metodo \texttt{summarize} riceve il
testo estratto da una pagina web e la domanda dell'utente, costruisce un
\textit{prompt} di sistema che impone regole stringenti e restituisce il risultato, 
oppure \texttt{None} se il contenuto non è
rilevante.

Entrambi gli adapter sono implementati come \textit{dataclass} immutabili
(\texttt{@dataclass(frozen=True)}), con le dipendenze necessarie
(\texttt{CacheRepository}, \texttt{OpenAIRepository}) iniettate direttamente
come campi della classe. Questa scelta, ricorrente in tutti gli adapter del
sistema, mantiene gli adapter privi di stato mutabile interno, semplificandone
test e sostituibilità.

\begin{figure}[H]
    \centering
    \includegraphics[width=0.58\textwidth]{uml/adapter.png}
    \caption{Diagramma UML del pattern Adapter implementato per \texttt{CacheAdapter}}
\end{figure}
% \section{Codifica}
